\documentclass[11pt]{article}
\usepackage[margin=1in]{geometry}
\usepackage{amsmath,amssymb,amsthm}
\usepackage{setspace}
\usepackage[hidelinks]{hyperref}
\hypersetup{hidelinks,colorlinks=false,pdfborder={0 0 0}}
\onehalfspacing
\newtheorem{proposition}{Proposition}
\newtheorem{lemma}{Lemma}
\newtheorem{example}{Example}
\newcommand{\E}{\mathbb E}
\newcommand{\I}{\mathbf 1}
\title{Allocating Heterogeneous Public Slots with Hidden Welfare Weights}
\author{Anonymous manuscript}
\date{October 2026}
\begin{document}
\maketitle

\begin{abstract}
How should a designer assign scarce slots of unequal quality when willingness to pay and individual social welfare weights are unobserved, group labels are observed, and transfers to recipients are restricted? I study independent private values and latent welfare weights under dominant-strategy incentives, ex post participation, and no positive transfers. A report of a welfare weight that does not enter private utility cannot generally reveal that weight. The relevant statistic is instead its conditional expectation given an incentive-compatible value report and an observable group. The welfare objective then has a virtual-surplus representation that accounts for both informational rents and the public value of collected charges. When each group's virtual score rises with value, assigning qualities in descending score order, with envelope payments, is optimal within the stated mechanism class. A bounded-subsidy extension leaves this allocation rule intact when transfers have no common budget. An explicit irregular example shows why pointwise score ranking otherwise fails incentive compatibility; a second example shows that free, report-independent assignment can be optimal when revenue has no social value. These results describe a conditional design rule, not an unconditional ranking of groups or a policy recommendation.
\end{abstract}

\section{Introduction}

Publicly allocated places, permits, or service slots can be scarce and differ in quality. Charges may help elicit who values a place, yet the designer may value an extra dollar of public revenue differently from a recipient's dollar of utility. Moreover, the recipient's social welfare weight need not be observable even when a group label is. The allocation problem therefore combines private willingness to pay, hidden normative information, heterogeneous objects, and a constraint on transfers to people.

This paper gives a precise answer for a narrow, assignable-resource model. It distinguishes a welfare weight from a preference parameter: a person's weight enters the designer's objective but not that person's utility. A direct request for the weight is consequently uninformative under dominant-strategy incentives, apart from changes at exceptional values. Under a known prior, the designer can use the conditional mean weight associated with an incentive-compatible value report and the observed group. A payment identity turns the objective into an expected sum of quality times a group-specific virtual score. If those scores are nondecreasing, a simple assignment rule implements the pointwise optimum. If they fall, the same ranking can ask a person to lose a slot by reporting a higher value, and so cannot be implemented truthfully.

The classic auction problem with privately known willingness to pay is a useful point of reference; the verified publisher description of Myerson (1981) concerns one object and a seller's expected utility. The verified publisher description of Saez and Stantcheva (2016) concerns the use of social marginal welfare weights to aggregate gains and losses in tax reform. Those descriptions motivate the two ingredients of this model but do not establish its results. The derivations here are self-contained. The available source evidence does not establish where these particular weighted-quality results sit in the full literature, so no priority or novelty claim is made.

\section{Environment and mechanism class}\label{sec:model}

There are $n$ agents and $m\le n$ indivisible, rival slots. Slot $j$ has publicly known quality $q_j$, with
\[
q_1\ge q_2\ge\cdots\ge q_m>0.
\]
An agent receives at most one slot; a slot is assigned to at most one agent. Unassigned agents have quality zero. Agents' group labels $g_i$ are observable. The declared private-type domain in group $g$ is the rectangle $[0,\bar v_g]\times\mathcal W_g$, where $\mathcal W_g\subseteq[0,\infty)$ is a nonempty set of possible weights. Incentive constraints apply to every pair in this domain, including pairs assigned zero probability by the prior. Conditional on the vector of labels, the pairs $(v_i,w_i)$ are independent across agents. The value $v_i$ has an atomless distribution on $[0,\bar v_{g_i}]$, with strictly positive density $f_{g_i}$ almost everywhere and distribution function $F_{g_i}$. The social weight $w_i\ge0$ has finite expectation. The joint distribution of value and weight is known to the designer; their within-agent dependence is unrestricted. Write
\[
\mu_g(v)=\E[w_i\mid v_i=v,g_i=g],\qquad
\bar\mu_g=\E[w_i\mid g_i=g].
\]
Conditional means are defined almost everywhere. No assumption says that a group has one common welfare weight.

A reported profile determines a lottery over feasible matchings and a charge $p_i$ for each agent, where $p_i<0$ is a transfer \emph{to} the agent. Let $Q_i$ be the expected quality from the lottery. Quasilinear utility is
\begin{equation}\label{eq:utility}
u_i=v_iQ_i-p_i.
\end{equation}
The core mechanism class imposes dominant-strategy incentive compatibility (DSIC), ex post individual rationality (IR), and no positive transfers, $p_i\ge0$. These conditions hold for every report by opponents and every type of the agent, taking expectations only over a mechanism's lottery. They allow charges from recipients. Group labels are verified; $v_i$ and $w_i$ are not. The report domain for $(v_i,w_i)$ is rectangular, so a person can claim any allowed weight at any allowed value. A randomized mechanism's utility and welfare depend on its expected qualities and expected charges.

For a public-funds value $\lambda\ge0$, the designer maximizes
\begin{equation}\label{eq:objective}
W=\E\!\left[\sum_{i=1}^n w_i(v_iQ_i-p_i)+\lambda\sum_{i=1}^n p_i\right].
\end{equation}
Thus $w_i$ weights the recipient's \emph{net} utility and $\lambda$ weights revenue available to the public budget. This is a normative objective, not an identified empirical welfare measure. The known joint prior and the specified $\lambda$ are maintained inputs. If either is unknown to the designer, the optimization below cannot be implemented as stated.

The word ``public'' here describes the provider. Slots are rival and assignable, not a nonrival public good. The transfer restriction is a ban on payments \emph{to} agents; a ceiling on charges \emph{from} agents is a different problem.

\section{What private weights can reveal}\label{sec:weights}

Because $w_i$ has no effect on utility in \eqref{eq:utility}, reporting it cannot generally supply the missing individual welfare weight.

\begin{proposition}[Weight-report equivalence]\label{prop:weights}
Fix an agent and opponents' reports. In any DSIC mechanism, truthful utility is the same for all reported weights at a given true value. At every value where that utility is differentiable, truthful expected quality and charge are also the same for every weight report. Consequently, with an atomless value distribution, weight reports cannot change expected qualities or charges on a positive-probability set of true values.
\end{proposition}

The complete proof is in Appendix \ref{app:proofs}. At a kink of the utility function, distinct weight reports could select distinct subgradients, so the proposition is an almost-everywhere statement. It also does not rule out elicitation when a weight affects preferences, is verifiable, or has some other contractual consequence. In the present model, the welfare-relevant statistic after a truthful value report is $\mu_{g_i}(v_i)$.

\section{Payments and virtual welfare}\label{sec:score}

Suppress opponents' reports and a weight report. DSIC for the scalar value implies a nondecreasing quality function $Q_i(v)$ and the envelope identity
\begin{equation}\label{eq:envelope}
U_i(v)=U_i(0)+\int_0^v Q_i(t)\,dt,
\qquad
p_i(v)=vQ_i(v)-U_i(0)-\int_0^vQ_i(t)\,dt.
\end{equation}
Conversely, a nondecreasing $Q_i$ and payments given by \eqref{eq:envelope} are DSIC. At $v=0$, IR and $p_i\ge0$ imply $U_i(0)=p_i(0)=0$. Monotonicity then implies $p_i(v)\ge0$ at all values, and \eqref{eq:envelope} implies IR.

Define the group score, at values with $f_g(v)>0$, by
\begin{equation}\label{eq:score}
S_g(v)=\lambda v-\lambda\frac{1-F_g(v)}{f_g(v)}
+\frac{\int_v^{\bar v_g}\mu_g(t)f_g(t)\,dt}{f_g(v)}.
\end{equation}
It combines the public value of payments with the value of preserving recipient rents. The last term uses an \emph{upper-tail integral} of social weights, not just $\mu_g(v)$: changing quality at a low value changes payments and utility at higher values through incentive constraints.

\begin{proposition}[Virtual-welfare identity]\label{prop:identity}
Under the assumptions of Section \ref{sec:model}, every feasible DSIC, ex post IR mechanism without positive transfers satisfies
\begin{equation}\label{eq:virtual}
W=\E\!\left[\sum_{i=1}^n S_{g_i}(v_i)Q_i(\mathbf v)\right],
\end{equation}
where the equality is unchanged by the reports of weights except on value-null sets. The qualities in \eqref{eq:virtual} may be randomized.
\end{proposition}

The proof in Appendix \ref{app:proofs} substitutes the payment identity into \eqref{eq:objective} and reverses a bounded double integral. Cross-agent independence ensures that conditioning on opponents' reports leaves $f_g$ and $\mu_g$ unchanged. Equation \eqref{eq:virtual} is an accounting identity, not yet an implementation theorem: a pointwise maximizer of its right-hand side need not be monotone in an agent's own value.

\section{A regular assignment rule}\label{sec:regular}

Call a group \emph{regular for this objective} if $S_g$ has a nondecreasing version on its value support. This depends on the prior, the conditional social weights, and $\lambda$; regularity is a substantive assumption.

\begin{proposition}[Optimal assignment under regularity]\label{prop:optimal}
Suppose each observed group is regular. At each value-report profile, order agents by decreasing $S_{g_i}(v_i)$. Assign the best slot to the highest positive score, the next-best slot to the next positive score, and continue until either slots or positive scores run out. Leave unused slots unassigned. Resolve equal scores by a fixed priority independent of reports. For each agent, set the charge using \eqref{eq:envelope} with $U_i(0)=0$. This mechanism is feasible, DSIC, ex post IR, has no positive transfers, and maximizes \eqref{eq:objective} among mechanisms with those properties.
\end{proposition}

The proof is in Appendix \ref{app:proofs}. Sorting maximizes virtual welfare at each profile by pairwise swaps. Regularity makes an agent's assigned quality nondecreasing in that agent's value, including across fixed-priority ties. The envelope charge therefore implements the assignment. The virtual-welfare identity then gives a global upper bound attained by this mechanism. The rule can deliberately leave a positive-quality slot unused when all remaining virtual scores are negative.

For a transparent special case, let $v\sim U[0,1]$ and let a group's weight be the constant $a_g$. Direct substitution gives
\begin{equation}\label{eq:linear}
S_g(v)=(2\lambda-a_g)v+(a_g-\lambda).
\end{equation}
For $\lambda>0$, the score is nondecreasing exactly when $a_g\le2\lambda$. With $a_g=0$ it crosses zero at $v=1/2$, reproducing a revenue-only reserve in the one-slot case. With $a_g=\lambda$, $S_g(v)=\lambda v$, so allocation follows value. With $a_g=2\lambda$, the score is flat: value reports do not affect the virtual-welfare ranking within that group, and fixed-priority ties are truthful. These are model implications, not estimates of actual group weights.

\subsection*{A bound on transfers to recipients}

The no-transfer core sets the permitted payment to an individual to zero. Suppose instead that each agent may receive at most $B_i\ge0$, so $p_i\ge-B_i$, with no common budget constraint. The same envelope identity applies. At value zero, IR and the transfer bound give $0\le U_i(0)\le B_i$. For a nondecreasing allocation, this condition is also sufficient for IR and the transfer bound at every value. Repeating the objective calculation adds $(\bar\mu_{g_i}-\lambda)U_i(0)$ to that agent's expected virtual welfare, conditional on opponents' reports. Thus, for each agent, an optimal mechanism sets $U_i(0)=B_i$ if $\bar\mu_{g_i}>\lambda$, sets it to zero if $\bar\mu_{g_i}<\lambda$, and may choose any value in the interval at equality. This is a fixed subsidy that does not change allocation incentives. Under score regularity, Proposition \ref{prop:optimal}'s assignment remains optimal with this modified payment. Appendix \ref{app:subsidy} verifies the extension. A shared transfer budget would couple agents' choices and is outside this result.

\section{When ranking fails}\label{sec:limits}

\begin{example}[A falling score violates incentives]\label{ex:fall}
Take two agents, one unit-quality slot, independent uniform values on $[0,1]$, constant weights $w_i=3$, and $\lambda=1$. Equation \eqref{eq:linear} gives $S(v)=2-v$. Hold agent 2's report at $1/2$. The pointwise score maximizer assigns the slot to agent 1 when she reports $1/4$, but assigns it to agent 2 when she reports $3/4$. If agent 1's true value moves from $1/4$ to $3/4$, her assigned quality falls from one to zero. The two incentive inequalities for these values require $(3/4-1/4)(Q(3/4)-Q(1/4))\ge0$, which is false. Thus raw-score ranking is not DSIC and Proposition \ref{prop:optimal} does not apply.
\end{example}

An optimum for falling scores must optimize over monotone allocations, for instance through an appropriate ironing or direct incentive-constrained program. This manuscript gives no general ironing theorem for heterogeneous groups. The failure above is already enough to reject an unconditional recommendation to rank by \eqref{eq:score}.

\begin{proposition}[A free-assignment benchmark]\label{prop:free}
Suppose there are two agents, one unit-quality slot, independent uniform values on $[0,1]$, constant weights $w_i=1$, and $\lambda=0$. Every feasible DSIC, ex post IR, no-positive-transfer mechanism has $W\le1/2$. Assigning the slot to either fixed agent for free, or by a value-independent lottery for free, attains $W=1/2$.
\end{proposition}

The proof in Appendix \ref{app:proofs} uses $S(v)=1-v$ and the negative covariance of a decreasing score with a nondecreasing own-value allocation. Charging for a value-sensitive assignment destroys recipient utility when revenue has no social value. The benchmark does not say that free assignment is optimal for other distributions, group weights, or values of $\lambda$.

\section{Discussion and limitations}\label{sec:discussion}

The regular rule answers the design question only after specifying a joint prior for $(v,w)$, the public value $\lambda$, and the exact transfer restriction. Even when $w$ itself is hidden, a known correlation between $w$ and $v$ can affect $\mu_g(v)$ and thus the score. Without a defensible prior or welfare criterion, group-specific scores cannot be calculated. The analysis offers no empirical method to identify either object.

The model also assumes independent agents conditional on group labels, one-dimensional multiplicative values for quality, quasilinearity, verifiable group labels, and no externalities among recipients. It imposes ex post IR and DSIC, which can be stronger than participation and incentive requirements in particular institutions. A cap on \emph{charges from} agents, $p_i\le C_i$, would alter implementability and potentially the optimal allocation. A genuine nonrival public good, complementary qualities, or fairness constraints needs a different feasibility and welfare analysis. These boundaries matter more than the tractability of the score formula.

The limited verified literature record supports background links to the private-value auction problem and to welfare-weighted tax analysis. It does not justify a claim that the propositions are new or that they resolve an existing frontier. The paper's claim scope is theoretical: under the stated assumptions and regularity, the rank-by-score mechanism is optimal in the stated class; outside regularity, the examples show why a different analysis is required.

\section{Conclusion}

For rival public slots of unequal quality, latent welfare weights cannot be elicited by a costless weight report when they do not enter utility. With a known prior, observable group and truthful value report supply a conditional mean weight. The envelope formula converts weighted net welfare plus public revenue into virtual welfare. Monotone group scores permit quality assignment by score and incentive-compatible charges without transfers to individuals. Falling scores or a different transfer rule call for a separate constrained optimization. The conclusions remain conditional on the model's information, welfare, and feasibility assumptions.

\appendix
\section{Complete paper proofs}\label{app:proofs}

\subsection{Proposition \ref{prop:weights}: weight-report equivalence}

Fix agent $i$ and every report of the other agents. Write $(Q(r,z),p(r,z))$ for the expected quality and charge offered by value report $r$ and weight report $z$. Utility from this report for true value $v$ is $vQ(r,z)-p(r,z)$, independent of the true weight. For any two true weights $z,z'$ and the same $v$, DSIC gives both
\[
vQ(v,z)-p(v,z)\ge vQ(v,z')-p(v,z')
\]
and the reverse inequality. Thus the truthful utility is a common number $U(v)$ for every weight. DSIC against every possible $(r,z)$ also gives
\[
U(v)=\sup_{r,z}\{vQ(r,z)-p(r,z)\}.
\]
As a supremum of affine functions of $v$, $U$ is convex. Bounded qualities in $[0,q_1]$ make it $q_1$-Lipschitz: an agent can compare the maximizing contracts at two values to bound the utility difference by $q_1$ times the value difference. The truthful contract at $(v,z)$ supports $U$ at $v$, since for every $t$ DSIC implies
\[
U(t)\ge tQ(v,z)-p(v,z)=U(v)+(t-v)Q(v,z).
\]
Hence $Q(v,z)$ is a subgradient of $U$ at $v$. A convex Lipschitz function on an interval is differentiable almost everywhere; at such a point the only subgradient is $U'(v)$. Therefore $Q(v,z)=U'(v)$ for all weight reports $z$, and $p(v,z)=vU'(v)-U(v)$ for all $z$. Endpoints and nondifferentiability points have zero probability under the atomless value law. This proves the stated almost-everywhere equivalence. The argument concerns expected quality and charge; lotteries with the same expectations need not be identical.

\subsection{Envelope lemma and Proposition \ref{prop:identity}}

Fix opponents' reports. For any two values $s<t$, the two value IC inequalities yield
\[
sQ(s)-p(s)\ge sQ(t)-p(t),\qquad
tQ(t)-p(t)\ge tQ(s)-p(s).
\]
Adding them gives $(t-s)(Q(t)-Q(s))\ge0$, so $Q$ is nondecreasing. The same IC inequalities imply that $Q(v)$ is a subgradient of the truthful utility $U(v)$, which is convex and Lipschitz by the argument above. Its derivative equals $Q(v)$ almost everywhere, and absolute continuity gives $U(v)=U(0)+\int_0^vQ(t)dt$. Rearranging $U(v)=vQ(v)-p(v)$ proves \eqref{eq:envelope}.

Conversely, let $Q$ be nondecreasing and define $U$ and $p$ by \eqref{eq:envelope}. For true $v$ and report $r$, truthful utility minus reporting utility is
\[
U(v)-\{vQ(r)-p(r)\}
=\int_r^v\{Q(t)-Q(r)\}\,dt\ge0.
\]
When $v<r$, the integral reverses its limits and its integrand is nonpositive, so the result still holds. At $v=0$, $U(0)=-p(0)$; IR requires $U(0)\ge0$, while no positive transfers requires $p(0)\ge0$. Hence both equal zero. Since $Q(t)\le Q(v)$ for $t\le v$, $p(v)=vQ(v)-\int_0^vQ(t)dt\ge0$. And $U(v)=\int_0^vQ(t)dt\ge0$ because qualities are nonnegative. These facts include $v=0$ and any flat or jump regions of $Q$.

To derive the virtual identity, fix the opponents' reports and a group $g$. Independence gives the conditional law of $(v,w)$ as $f_g$ and $\mu_g$ even after those reports. Suppress indices and set $b=\bar v_g$. Conditional expected welfare for this agent is
\begin{align*}
&\int_0^b f_g(v)\bigl[\mu_g(v)vQ(v)
+(\lambda-\mu_g(v))p(v)\bigr]dv\\
&\quad=\int_0^b f_g(v)\lambda vQ(v)\,dv
-\int_0^b f_g(v)(\lambda-\mu_g(v))
\left[\int_0^vQ(t)dt\right]dv\\
&\quad=\int_0^b\left\{\lambda vf_g(v)
-\int_v^b(\lambda-\mu_g(s))f_g(s)ds\right\}Q(v)dv\\
&\quad=\int_0^b f_g(v)S_g(v)Q(v)dv.
\end{align*}
The second line uses $U(0)=0$ and \eqref{eq:envelope}. The third reverses the order of integration over $0\le t\le v\le b$; this is valid because $Q\le q_1$, $b<\infty$, and $\E[w\mid g]<\infty$, so the absolute double integral is finite. The last line uses $\int_v^b f_g(s)ds=1-F_g(v)$ and definition \eqref{eq:score}. Sum over agents and integrate over opponents' reports. Proposition \ref{prop:weights} makes possible weight-report distinctions null for this expectation. This proves Proposition \ref{prop:identity}.

\subsection{Proposition \ref{prop:optimal}: regular assignment}

At a fixed value profile let $s_i=S_{g_i}(v_i)$. Consider any feasible deterministic matching. If agent $a$ has score $s_a\ge s_b$ but receives quality $q_l<q_h$ received by agent $b$, swapping their slots changes virtual surplus by $(s_a-s_b)(q_h-q_l)\ge0$. The same calculation applies if one quality is zero. Assigning a positive-quality slot to a nonpositive score cannot increase virtual surplus relative to leaving it unused. Repeated pairwise swaps therefore produce the stated descending-score assignment and show that it maximizes $\sum_i s_iQ_i$ among all deterministic matchings. A lottery is a convex combination of matchings, so its expected virtual surplus cannot exceed that maximum. Equality cases are harmless; fixed priority chooses one maximizing matching.

Fix an agent and all other reports. When that agent raises her value, regularity makes her score weakly rise. Her position in the score ordering can only improve, and fixed tie priority cannot make her position worse. She therefore receives weakly higher quality, including when she moves from no assignment to a slot. Her $Q(v)$ is nondecreasing. The envelope lemma then proves DSIC, IR, and nonnegative charges for the formula in \eqref{eq:envelope}. It also shows feasibility is unaffected by payments. Finally, Proposition \ref{prop:identity} bounds the welfare of every admissible mechanism by the expected pointwise maximum of virtual surplus. The constructed mechanism attains that maximum at every profile, hence attains the bound. This proves optimality without assuming that every pointwise maximizing rule is implementable; regularity supplies the needed monotonicity here.

\subsection{Bounded-subsidy extension}\label{app:subsidy}

Keep a fixed opponent-report profile and allow $p(v)\ge-B_i$. Equation \eqref{eq:envelope} gives $p(0)=-U(0)$, so IR at zero and the transfer bound imply $0\le U(0)\le B_i$. Conversely, if this interval condition holds and $Q$ is nondecreasing, then $U(v)=U(0)+\int_0^vQ(t)dt\ge0$ and
\[
p(v)=\underbrace{vQ(v)-\int_0^vQ(t)dt}_{\ge0}-U(0)\ge-B_i.
\]
In the conditional-welfare derivation for Proposition \ref{prop:identity}, retaining $U(0)$ adds
\[
-U(0)\int_0^{\bar v_g}(\lambda-\mu_g(v))f_g(v)dv
=(\bar\mu_g-\lambda)U(0).
\]
The coefficient depends only on the group, by cross-agent independence; $U(0)$ may be chosen separately for each opponent-report profile. Maximizing this linear term over $[0,B_i]$ gives the stated endpoint choices, including any choice at equality. These choices shift all payments for the agent by the same amount and therefore preserve value IC. The virtual-surplus part still has the same pointwise upper bound and, under regularity, the same maximizing assignment. This proves the extension. If a common budget links transfers, the independent endpoint maximizations no longer follow.

\subsection{Proposition \ref{prop:free}: a no-revenue benchmark}

With $\lambda=0$, $\mu(v)=1$, and uniform $f(v)=1$, equation \eqref{eq:score} gives $S(v)=1-v$. Fix an agent and an opponent's report. Let $V,V'$ be independent uniform draws and $Q$ the agent's nondecreasing own-value allocation rule. The product
\[
\bigl[(1-V)-(1-V')\bigr]\bigl[Q(V)-Q(V')\bigr]\le0
\]
for every pair of draws. Expanding its expectation and using independence gives
\[
2\,\operatorname{Cov}(1-V,Q(V))\le0.
\]
Because $\E[1-V]=1/2$, it follows that $\E[(1-V)Q(V)]\le\tfrac12\E[Q(V)]$. Integrate over the opponent's report, apply this bound to both agents, and use feasibility $Q_1+Q_2\le1$ at every profile. Proposition \ref{prop:identity} then yields
\[
W\le\tfrac12\E[Q_1+Q_2]\le\tfrac12.
\]
A fixed assignment of the unit slot for zero charge is DSIC, IR, and transfer-free; its recipient's expected value is $1/2$, so its welfare equals the bound. Any report-independent lottery that always assigns the slot for free has the same welfare. This proves Proposition \ref{prop:free}.

\clearpage
\subsection{Proof-obligation record}

The weight-report argument uses the elementary fact that convex Lipschitz functions on an interval are differentiable almost everywhere and equal the integral of their almost-everywhere derivative; the hypotheses are established above by bounded quality. The virtual identity's interchange is justified by finite value bounds, bounded quality, and finite expected weight. The ranking proof covers ties, unassigned slots, and lotteries. No theorem in this manuscript is claimed to be machine verified. The optimal allocation for nonmonotone scores is unresolved here.

\begin{thebibliography}{9}
\bibitem{Myerson1981}
Myerson, Roger B. 1981. ``Optimal Auction Design.'' \emph{Mathematics of Operations Research} 6 (1): 58--73. \url{https://doi.org/10.1287/moor.6.1.58}.
\bibitem{SaezStantcheva2016}
Saez, Emmanuel, and Stefanie Stantcheva. 2016. ``Generalized Social Marginal Welfare Weights for Optimal Tax Theory.'' \emph{American Economic Review} 106 (1): 24--45. \url{https://doi.org/10.1257/aer.20141362}.
\end{thebibliography}

\end{document}
